\documentclass[11pt]{article}
\usepackage[a4paper,margin=28mm]{geometry}
\usepackage{amsmath,amssymb,amsthm,booktabs}
\usepackage[hidelinks]{hyperref}
\newtheorem{theorem}{Theorem}
\newcommand{\RA}{\operatorname{RA}}
\newcommand{\Prob}{\mathbb P}
\hypersetup{pdftitle={A three-leaf counterexample for Fitch reconstruction under a molecular clock},pdfauthor={Kaiwen Jin}}
\title{A three-leaf counterexample for Fitch reconstruction\protect\\ under a molecular clock}
\author{Kaiwen Jin\thanks{\href{mailto:yui.ui.siki@gmail.com}{yui.ui.siki@gmail.com}}}
\date{October 15, 2026}
\begin{document}
\maketitle
\begin{abstract}
Does Fitch parsimony on a tree whose leaves are all the same distance from the
root reconstruct the root state at least as accurately as any single leaf?
We give a negative answer for the symmetric twenty-state model. A rooted binary
tree with three leaves and strictly positive edge lengths has Fitch accuracy
$38829/400000$, whereas each leaf agrees with the root with probability
$39/400$. The strict deficit is $171/400000$. The proof is an exact calculation
over four types of Fitch set and includes uniform final tie breaking.
\end{abstract}

\section{Model and result}

Ancestral-state reconstruction estimates a character at a common ancestor from
its observed values at present-day leaves. Li, Steel and Zhang
\cite[p.~652]{LSZ08} conjectured that, under a symmetric model on an ultrametric
tree, Fitch parsimony using all leaves is at least as accurate as using any one
leaf. Here ultrametric means that every root-to-leaf path has the same total
length. The two-state case was proved by Fischer and Thatte
\cite[author version, Theorem~2]{FT09}, and the three- and four-state cases by Herbst and Fischer
\cite[author version, Theorems~1--2]{HF18}. We exhibit a twenty-state counterexample. It negates
the universal assertion over alphabet sizes $r\geq5$, without asserting failure
at each such size.

Let $r\geq2$ be an integer. On the alphabet $\{1,\ldots,r\}$, choose the
state at the root $\rho$ uniformly, and evolve states independently down
branches conditional on their parent states. Write $\sigma_x$ for the state
at a vertex $x$. An edge
of length $\ell\geq0$ uses the continuous-time Markov generator
\begin{equation}\label{eq:generator}
 Q_{ij}=\begin{cases}-1,&i=j,\\1/(r-1),&i\ne j.\end{cases}
\end{equation}
Thus the expected substitution rate is one. If every leaf has height $H$,
the probability that a specified leaf has the root state is
\begin{equation}\label{eq:benchmark}
 q_r(H)=\frac1r+\frac{r-1}{r}\exp\!\left(-\frac{rH}{r-1}\right).
\end{equation}

The Fitch set at a leaf is its observed singleton. At each internal vertex,
take the intersection of the two child sets when it is nonempty and their union
otherwise. Choose the root estimate $\widehat\sigma_\rho$ uniformly from the
resulting root set. Write
$\RA=\Prob(\widehat\sigma_\rho=\sigma_\rho)$, including this final random choice.
Only the topology and leaf states enter the set recursion.

\begin{theorem}\label{thm:counterexample}
Let $r=20$. Let the root $\rho$ have children $a,v$, and let $v$ have children
$b,c$, so that the leaves are $a,b,c$. Give the edges the lengths
\begin{equation}\label{eq:lengths}
 \ell_{\rho a}=\frac{19}{20}\log20,\qquad
 \ell_{\rho v}=\frac{19}{20}\log16,\qquad
 \ell_{vb}=\ell_{vc}=\frac{19}{20}\log\frac54.
\end{equation}
All four lengths are finite and strictly positive, and every leaf has height
$H=(19/20)\log20$. Under \eqref{eq:generator} and the uniform root and tie-breaking
conventions above,
\begin{equation}\label{eq:result}
 \RA=\frac{38829}{400000}
   <\frac{39}{400}=q_{20}(H),\qquad
 q_{20}(H)-\RA=\frac{171}{400000}.
\end{equation}
\end{theorem}

\section{Exact calculation}

\begin{proof}[Proof of Theorem~\ref{thm:counterexample}]
All logarithms are natural. The clock condition follows from
$\log16+\log(5/4)=\log20$. To obtain the edge probabilities with the stated rate,
let $\Pi$ be the matrix all of whose entries are $1/r$. Since $\Pi^2=\Pi$ and
$Q=-r(I-\Pi)/(r-1)$, the exponential series gives
\begin{equation}\label{eq:channel}
 e^{\ell Q}=\Pi+\lambda(I-\Pi),\qquad
 \lambda=\exp\!\left(-\frac{r\ell}{r-1}\right).
\end{equation}
In particular, an edge preserves its input state with probability
$(1+(r-1)\lambda)/r$ and changes it to each specified different state with
probability $(1-\lambda)/r$.

Condition for now on $\sigma_\rho=1$. The values of $\lambda$ on the edges
$\rho a$, $\rho v$, $vb$ and $vc$ are $1/20$, $1/16$, $4/5$ and $4/5$,
respectively. Hence
\begin{equation}\label{eq:primitive}
 \begin{aligned}
 q:=\Prob(\sigma_a=1\mid\sigma_\rho=1)&=\frac{39}{400},&
 \Prob(\sigma_a=j\mid\sigma_\rho=1)&=\frac{19}{400}=:p\quad(j\ne1),\\
 \Prob(\sigma_v=1\mid\sigma_\rho=1)&=\frac7{64},&
 \Prob(\sigma_v\ne1\mid\sigma_\rho=1)&=\frac{57}{64}.
 \end{aligned}
\end{equation}
Given any state at $v$, each of $b,c$ equals that state with probability
$u=81/100$ and equals each specified different state with probability
$w=1/100$. The two terminal evolutions are independent conditional on
$\sigma_v$.

Let $B$ be the Fitch set at $v$. It is a singleton when $\sigma_b=\sigma_c$
and a two-element set otherwise. Its total singleton probability is
$u^2+19w^2=329/500$, independent of the state at $v$. Denote the probabilities
of the four possibilities by $c_1,n_1,c_2,n_2$, where $c_k$ corresponds to
$|B|=k$ and $1\in B$, and $n_k$ to $|B|=k$ and $1\notin B$. All probabilities
in the next calculation remain conditional on $\sigma_\rho=1$. We have
\begin{equation}\label{eq:cherry}
 \begin{aligned}
 c_1&=\frac7{64}u^2+\frac{57}{64}w^2=\frac{1437}{20000},&
 n_1&=\frac{329}{500}-c_1=\frac{11723}{20000},\\
 c_2&=\frac7{64}\,2u(1-u)+\frac{57}{64}\,2w(1-w)
       =\frac{513}{10000},&
 n_2&=1-\frac{329}{500}-c_2=\frac{2907}{10000}.
 \end{aligned}
\end{equation}
Indeed, $c_1$ is the probability that both terminal states equal 1, and $c_2$
is the probability that exactly one of them equals 1. These four events are
disjoint, exhaustive and have probabilities summing to one.

Conditional on the root, the leaf $a$ is independent of the entire subtree
below $v$, and therefore independent of $B$. Fix a possible set $B$ of size
$k\in\{1,2\}$. If $1\in B$, the value $\sigma_a=1$ produces a correct singleton
at the root. A value in $B\setminus\{1\}$ produces an incorrect singleton.
Each of the $20-k$ states outside $B$ produces the union $B\cup\{\sigma_a\}$
of size $k+1$, with success probability $1/(k+1)$. Thus the conditional
success probability is
\begin{equation}\label{eq:payoffs}
 \begin{cases}
 q+(20-k)p/(k+1),&1\in B,\\
 q/(k+1),&1\notin B.
 \end{cases}
\end{equation}
For the second case, only $\sigma_a=1$ can yield success, and it yields a
union of size $k+1$. No state is randomly selected at $v$; the full set $B$
is used in this merge.

Substituting $k=1,2$ into \eqref{eq:payoffs}, then averaging the four events
in \eqref{eq:cherry}, gives
\begin{align}
 \Prob(\widehat\sigma_\rho=1\mid\sigma_\rho=1)
 &=c_1\frac{439}{800}+n_1\frac{39}{800}
   +c_2\frac{153}{400}+n_2\frac{13}{400}\notag\\
 &=\frac{630843+457197+313956+151164}{16000000}
  =\frac{38829}{400000}.\label{eq:accuracy}
\end{align}
Permuting all state labels preserves the transition probabilities, the Fitch
set operations and the uniform final choice. The same conditional accuracy
therefore holds for every root state, and averaging the uniform root law
gives the asserted value of $\RA$.

Finally, \eqref{eq:benchmark} and $\exp(-20H/19)=1/20$ give
$q_{20}(H)=39/400=39000/400000$. Subtraction proves the strict deficit in
\eqref{eq:result}. The tree and lengths in \eqref{eq:lengths} satisfy every
premise of the proposed universal inequality, so this single example disproves
that assertion.
\end{proof}

\paragraph{Verification material.}
The accompanying script \texttt{code/exact\_three\_leaf.py} independently
sums over all $20^4$ assignments to the four non-root vertices and applies the
literal Fitch recursion. From the package root, run
\texttt{python3 code/exact\_three\_leaf.py} with Python 3.10 or newer; only
the standard library is needed. The displayed finite calculation is the proof;
the script provides a reproducible check.

\paragraph{AI assistance.}
An AI assistant was used to develop and check the construction, prepare the
verification script, and draft the text. Two fresh-context AI reviews, one
focused on reading and one on mathematical objections, informed a single
revision. Independent mathematical review is requested.

\begin{thebibliography}{9}
\bibitem{LSZ08}
G.~Li, M.~Steel and L.~Zhang,
More Taxa Are Not Necessarily Better for the Reconstruction of Ancestral
Character States,
\emph{Systematic Biology} \textbf{57}(4) (2008), 647--653.
\href{https://doi.org/10.1080/10635150802203898}{doi:10.1080/10635150802203898}.

\bibitem{FT09}
M.~Fischer and B.~D.~Thatte,
Maximum Parsimony on Subsets of Taxa,
\emph{Journal of Theoretical Biology} \textbf{260}(2) (2009), 290--293.
\href{https://doi.org/10.1016/j.jtbi.2009.06.010}{doi:10.1016/j.jtbi.2009.06.010}.
Author manuscript: \href{https://arxiv.org/abs/0809.3653v2}{arXiv:0809.3653v2}.

\bibitem{HF18}
L.~Herbst and M.~Fischer,
On the Accuracy of Ancestral Sequence Reconstruction for Ultrametric Trees
with Parsimony,
\emph{Bulletin of Mathematical Biology} \textbf{80} (2018), 864--879.
\href{https://doi.org/10.1007/s11538-018-0407-5}{doi:10.1007/s11538-018-0407-5}.
Author manuscript: \href{https://arxiv.org/abs/1706.06085v1}{arXiv:1706.06085v1}.
\end{thebibliography}
\end{document}
